\section{Introduction}
\label{sec:intro}
A locally valid action can lead away from a target. Completing a folding sequence therefore
requires selecting among alternatives, tracking their consequences, and revising earlier
decisions. We study long-horizon planning when the model receives the exact coordinates,
orientations, stacking order and original-sheet regions of every layer, and every action is
checked. Origami provides a generated state space with programmatically checkable success;
we use it as a controlled planning task rather than a model of physical folding.

We introduce \textbf{CP2Seq}: 600 generated instances, a configurable generator, and an
interactive evaluation environment. Given a crease pattern and a target folded state, the
model constructs a sequence of simple folds from an unfolded square. Reference sequences
and intermediate reference states are withheld; any valid sequence producing both the
specified creases and the target stack is accepted. Original-sheet identity and layer order
distinguish geometrically overlapping constructions (Section~\ref{sec:scoring}).

Existing origami benchmarks assess visual questions, pattern prediction or generation
\citep{gamibench,origamispace,origamibench}. CP2Seq instead tests executable sequence
recovery under explicit action constraints. Fold lines are specified in current folded
coordinates while target creases are given on the original sheet; planning connects these
representations through reflections and layer connectivity. Tools expose legality checks,
backtracking, optional legal-action enumeration and target feedback. Search baselines use
the same candidate actions, allowing us to distinguish task reachability from model success.

Our contributions are (1) a generated dataset and generator spanning reference depths and
two simple-fold action models; (2) an interactive pipeline with step-level checks and
replay-based terminal scoring, including original-sheet identity; and (3) matched search
baselines, assistance and modality ablations, repeated trials, and observable trajectory
analysis. These isolate measurable behavior without attributing it to unobserved internal
capabilities.

The models solve some short sequences but no medium or hard ablation instance. Yet random
legal folding solves 57\% of easy and 96\% of some-layers episodes within 80 folds, and BFS
solves five of ten medium instances. Removing images does not reduce solve rate. Failures
rarely reflect reaching the target without stopping; models explore fewer distinct states
than random folding and revisit discarded branches. Repeated trials reveal substantial
variance, while a second model family solves 14/20 instances in one episode each. We
therefore frame CP2Seq as planning over explicit geometric state, not as a validated measure
of visual perception or an isolated internal reasoning faculty. Depth and layer count are
confounded in the released all-layers corpus, limiting depth-based conclusions.
